
% \begin{table*}[!ht]
% \centering
% \caption{Background and Contact Information of Workshop Organizers}
% \label{authorBackground}
% %
% \begin{tabular}{p{2.625cm}p{4.1cm}p{5.5cm}p{3.5cm}}
% \toprule
% Presenter & Education & Current Position & Email \\
% \midrule
% Tushar Athawale & Ph.D., University of Florida & Staff Scientist, ORNL & athawaletm@ornl.gov \\
% Chris R. Johnson & Ph.D., University of Utah & Distinguished Professor, University of Utah & crj@sci.utah.edu \\
% Kristi Potter & Ph.D., University of Utah & Senior Scientist, NREL & kristi.potter@nrel.gov \\
% Paul Rosen & Ph.D., Purdue University & Associate Professor, University of Utah & paul.rosen@utah.edu \\
% David Pugmire & Ph.D., University of Utah & Distinguished Scientist, ORNL & pugmire@ornl.gov \\
% Antigoni Georgiadou & Ph.D., Florida State University & Applied Mathematician, ORNL & georgiadoua@ornl.gov\\
% Tim Gerrits & Ph.D., Magdeburg University & Senior Scientist, RWTH Aachen University & gerrits@vis.rwth-aachen.de\\
% Nadia Boukhelifa & Ph.D., University of Kent  & Researcher, INRAE/Univ. Paris-Saclay & nadia.boukhelifa@inrae.fr\\
% %Robert Maynard & B.S., Laurentian University & Principal Engineer, Kitware, Inc. & robert.maynard@kitware.com \\
% \bottomrule
% \end{tabular}
% \vspace{-5mm}
% \end{table*}


%\textbf{uncertainty 2024 workshop:} October 13, 2024: 2:00 PM-5:00 PM AEDT (UTC+11)


\section{Qualifications and Background}

{\large Dr. Athawale} $\;$ is a Research Scientist in Visualization Group at the Oak Ridge National Laboratory (ORNL). He holds the position of Joint Faculty Assistant Professor in Electrical Engineering and Computer Science (EECS) Department of University of Tennessee, Knoxville. He has first-authored multiple publications foundational to progressing uncertainty visualization research~\cite{TA:Athawale:2021:topoMappingUncertaintyMarchingCubes, Athawale:2021:AJE,TA:2022:Athawale:UQMorseComplex,Athawale:2019:dbsUncertainPos} that appeared in IEEE VIS conferences and other top-tier journals. He served as the author of the uncertainty visualization section for the US Department of Energy's (DOE’s) 2022 visualization workshop report and was invited to the prestigious 2023 Dagstuhl workshop on uncertainty. 

{\large Dr. Johnson} $\;$ is a Distinguished Professor at the Scientific Computing and Imaging (SCI) Institute in University of Utah. He is an international leader and recipient of multiple awards in the field of Scientific Visualization. He is one of the first scientists who recognized the need for visualization of uncertainty~\cite{TA:Johnson:2003:nextStepVisErrors, TA:2004:Johnson:topScivis} and has co-authored with all workshop organizers on a wide range of uncertainty visualization techniques and applications. He organized IEEE VIS 2011 and Dagstuhl 2023 workshops on uncertainty visualization.

{\large Dr. Potter} $\;$ is a senior scientist and group manager of the data, analysis, and visualization group in the Computational Science Center at the National Renewable Energy Laboratory (NREL). She is an expert in uncertainty and ensemble visualization and has co-authored a number of publications with Johnson and Rosen~\cite{TA:Potter:2009:ensemble-vis, TA:Potter:2012:UQtaxonomy, muView} and organized uncertainty visualization workshops at IEEE VIS 2015 and most recently a 2023 Dagstuhl on the topic. 

{\large Dr. Rosen} $\;$ is an Associate Professor at the SCI Institute, University of Utah. He received his PhD from Purdue University in 2010. 
%Subsequently, he served as a Research Assistant Professor at SCI Institute and the School of Computing at the University of Utah until 2015. 
His broad research interests include geometry- and topology-based approaches with analysis of uncertainty for solving problems in scientific and information visualization~\cite{TA:Potter:2012:UQtaxonomy, muView}. He has received 6 Best Paper or Honorable Mention awards. He was the recipient of a NSF Early Career Award in 2019.

{\large Dr. Pugmire} $\;$ is a Distinguished Scientist and leader of Visualization Group at the ORNL. He holds the position of Joint Faculty Professor in EECS Department of University of Tennessee, Knoxville. He is expert in high-performance computing, particle tracking, and uncertainty visualization. He co-authored publications with Athawale and Johnson on techniques and software for efficient uncertainty visualization in tools like ParaView and VisIt~\cite{TA:2023:Wang:funmc2, TA:Athawale:2023:uncertainFibers}. 
% The background and contact information of all workshop organizers is provided in Table~\ref{authorBackground}.

Dr. Georgiadou is an Applied Mathematician at the Science Engagement Section, Algorithms \& Performance Analysis (APA) Group. She obtained her Ph.D. in Mathematics from Florida State University where she worked on optimization in stellar evolution applications. Dr. Georgiadou is the working group lead of the Uncertainty Quantification group at the Oak Ridge Leadership Computing Facility (OLCF)~\cite{ensLead}. She has served as the chair of the SC workshop on "Enabling Predictive Science with Optimization and Uncertainty Quantification in HPC" taking place at The International Conference for High Performance Computing, Networking, Storage, and Analysis (SC'23, SC'24).

{\large Dr. Gerrits} $\;$ is a senior scientist and group manager of the visualization group at RWTH Aachen University, Germany. His interests include (in-situ) visualization for HPC applications, immersive analysis, and uncertainty visualization for vector field ensembles and second-order tensor data~\cite{Gerrits:2019:EuroVis, schmitz24exploring}. He has been organizer of the IEEE SciVis Contest since 2023 and has received 2 best paper or honorable mention awards.

{\large Dr. Boukhelifa} $\;$ is a research scientist in visual analytics at INRAE, Université Paris-Saclay, France. In 2023, she co-organized a Dagstuhl seminar with Johnson and Potter and co-edited a CG\&A special issue on the topic~\cite{boukhelifa_et_al:DagRep.12.8.1,boukhelifa2023visualization,wood2012sketchy}. She has co-authored multiple publications on the visual encoding and interpretation of uncertainty \cite{boukhelifa2023visualization,boukhelifa2012evaluating,boukhelifa2017data,boukhelifa2009uncertainty} and recently explored its impact on decision making in agronomy and urban planning \cite{courteille2024using}.
